\section{Assignment 2 视角：从 correctness 到 performance}

Lecture 6 对应的 systems assignment 不只是“写出能跑的 kernel”，而是要求你 benchmark 和 profile 实现。一个 kernel 的评价至少包含三层：

\begin{longtable}{p{0.20\textwidth}p{0.35\textwidth}p{0.35\textwidth}}
\toprule
\textbf{层级} & \textbf{问题} & \textbf{证据} \\
\midrule
Correctness & 输出是否和 PyTorch reference 接近？ & \texttt{torch.allclose}、单元测试、极端 shape。 \\
Performance & 是否比 baseline 快？随 shape 如何 scaling？ & CUDA event benchmark，多次 trial，warmup。 \\
Attribution & 时间花在哪些 kernels 和 memory 操作上？ & PyTorch profiler、Nsight、kernel 名、occupancy/throughput 指标。 \\
\bottomrule
\end{longtable}

\begin{importantbox}{不要跳过 correctness}
高性能错误 kernel 没有意义。Triton 的 mask、stride、边界条件、dtype accumulation 都可能造成 subtle bug。Assignment 的正确姿势是先让 reference 对齐，再 benchmark，再 profile，再优化。
\end{importantbox}

\begin{warningbox}{性能比较要公平}
比较 naive、builtin、compiled、Triton 版本时，必须使用相同输入 shape、相同 dtype、相同 device，并包含 warmup。否则测到的可能是编译时间、CPU launch overhead、allocator 行为或缓存状态，而不是 kernel 本身。
\end{warningbox}

\subsection{从本讲走向下一讲}

本讲所有优化仍在单 GPU 内部。下一讲进入多 GPU 后，同样的原则会扩大：HBM 往返变成 GPU 间通信，thread blocks 变成 ranks 和 collectives，kernel fusion 的“少搬数据”变成 all-reduce、reduce-scatter、all-gather 的通信账本。

\begin{knowledgebox}{统一心智模型}
单 GPU kernel 优化和多 GPU parallelism 都是在问：数据在哪里、谁需要它、什么时候移动、能不能少移动或多复用。Lecture 6 学的是这个问题在一个 GPU 内部的版本。
\end{knowledgebox}

